\section*{CHƯƠNG 1. MỞ ĐẦU}
\addcontentsline{toc}{section}{\numberline{}CHƯƠNG 1. MỞ ĐẦU}
\setcounter{section}{1}
\setcounter{subsection}{0}
\setcounter{figure}{0}
\setcounter{table}{0}

\subsection{Bối cảnh và tính cấp thiết của đề tài}
Trong kỷ nguyên số hóa hiện nay, công nghệ nhận diện và xác thực khuôn mặt (Face Verification/Recognition) đã trở thành một phần không thể thiếu trong nhiều lĩnh vực quan trọng như an ninh biên giới, giao dịch ngân hàng (e-KYC), quản lý nhân sự và bảo mật thiết bị cá nhân. Các thuật toán học sâu hiện đại, tiêu biểu là các kiến trúc được tích hợp trong thư viện DeepFace (như VGG-Face, FaceNet, ArcFace), đã đạt được độ chính xác vượt bậc trong các điều kiện chụp thông thường, thậm chí vượt qua cả khả năng nhận diện của con người trong một số tập thử nghiệm tiêu chuẩn.

Tuy nhiên, sự phát triển mạnh mẽ của ngành công nghiệp phẫu thuật thẩm mỹ (Plastic Surgery) toàn cầu đã đặt ra một thách thức sinh trắc học nghiêm trọng và chưa được giải quyết triệt để. Các can thiệp phẫu thuật thẩm mỹ hiện đại — bao gồm nâng mũi (rhinoplasty), gọt hàm (jaw shaving), cắt mí (blepharoplasty) hay độn cằm (chin augmentation) — trực tiếp làm thay đổi cấu trúc hình học phi tuyến tính và kết cấu mô mềm của khuôn mặt. Sự thay đổi diện mạo có chủ đích này làm phá vỡ các đặc trưng nhận diện toàn thể (holistic features) mà các hệ thống học sâu truyền thống phụ thuộc vào. Khi một người trải qua phẫu thuật thẩm mỹ, khoảng cách vector đặc trưng (facial embedding) giữa ảnh trước và sau phẫu thuật của họ thường bị kéo giãn ra rất lớn, vượt qua ngưỡng chấp nhận của hệ thống, dẫn đến việc hệ thống từ chối xác thực một cách sai lệch (tăng tỷ lệ False Rejection Rate - FRR).

Việc từ chối xác thực danh tính sau phẫu thuật thẩm mỹ không chỉ gây bất tiện lớn cho người dùng trong các hoạt động e-KYC hay check-in xuất nhập cảnh, mà còn tạo ra kẽ hở lớn cho việc giả mạo danh tính hoặc trốn tránh trách nhiệm pháp lý. Do đó, việc nghiên cứu một giải pháp cải tiến thuật toán so khớp khuôn mặt có khả năng kháng lại sự thay đổi do phẫu thuật thẩm mỹ (Plastic Surgery-Resistant Face Verification) bằng cách tối ưu hóa các phân vùng cục bộ bền vững trên khuôn mặt là vô cùng cấp thiết.

\subsection{Đặt vấn đề}
Hiện nay, bài toán so khớp khuôn mặt trước và sau phẫu thuật thẩm mỹ gặp phải các thách thức kỹ thuật cốt lõi sau:
\begin{itemize}
    \item \textbf{Sự thất bại của các bộ trích xuất đặc trưng toàn thể (Holistic Extraction):} Các mô hình so khớp hàng đầu như VGG-Face hay ArcFace nén toàn bộ thông tin khuôn mặt thành một vector đặc trưng duy nhất. Khi một bộ phận như mũi hoặc cằm bị phẫu thuật làm thay đổi lớn, sự thay đổi cục bộ này sẽ tác động mạnh lên toàn bộ vector đặc trưng toàn cục, khiến mô hình coi đây là hai người hoàn toàn khác nhau.
    \item \textbf{Sự khan hiếm của dữ liệu huấn luyện phẫu thuật thẩm mỹ:} Các mô hình nhận diện khuôn mặt phổ biến được huấn luyện trên các tập dữ liệu lớn chứa ảnh tự nhiên (như LFW, CASIA-WebFace). Các tập dữ liệu này hầu như không chứa các cặp ảnh trước/sau phẫu thuật thẩm mỹ của cùng một cá nhân, khiến mô hình không thể tự học cách bỏ qua các biến đổi phẫu thuật.
    \item \textbf{Yêu cầu về giải thuật thích ứng thời gian thực (Real-time Adaptive Logic):} Hệ thống cần có khả năng tự động phân tích và phát hiện ra vùng nào trên khuôn mặt đã bị chỉnh sửa phẫu thuật để phân bổ lại trọng số so khớp, ưu tiên các vùng cấu trúc xương bất biến (như vùng mắt và trán) mà không làm tăng đáng kể độ trễ tính toán của hệ thống.
\end{itemize}

\subsection{Hạn chế của thư viện DeepFace trong bài toán phẫu thuật thẩm mỹ}
Thư viện DeepFace (chạy trên các mô hình nền tảng như VGG-Face, FaceNet) hiện là một trong những framework so khớp khuôn mặt mã nguồn mở mạnh mẽ nhất. Tuy nhiên, khi ứng dụng trực tiếp vào bài toán xác thực sau phẫu thuật thẩm mỹ, DeepFace bộc lộ những hạn chế rõ rệt:
\begin{itemize}
    \item \textbf{Cơ chế ra quyết định cứng nhắc:} DeepFace chỉ thực hiện so sánh một cặp ảnh mặt đầy đủ và trả về một khoảng cách duy nhất (global distance). Cơ chế này hoàn toàn không có khả năng phân biệt giữa việc thay đổi diện mạo do phẫu thuật thẩm mỹ cục bộ và việc thay đổi do hai người khác nhau hoàn toàn.
    \item \textbf{Hiện tượng trôi lệch điểm mốc giải phẫu (Landmark Drift):} Khi khuôn mặt bị sưng tấy nặng nề sau phẫu thuật hoặc thay đổi hình dạng hàm, các bộ dò điểm mốc (như OpenCV hay RetinaFace) dễ bị xác định sai vị trí mắt, mũi, miệng. Lỗi này làm trôi lệch bước căn chỉnh khuôn mặt (facial alignment), khiến ảnh cắt khuôn mặt đưa vào CNN bị lệch góc, làm giảm nghiêm trọng độ chính xác của vector embedding.
    \item \textbf{Thiếu tính linh hoạt cục bộ (Local Non-invariance):} DeepFace coi trọng số của mọi vùng trên khuôn mặt là tương đương nhau trong quá trình trích xuất vector. Trong thực tế phẫu thuật thẩm mỹ, vùng mắt/trán thường có tính ổn định cao hơn rất nhiều so với vùng mũi và cằm, nhưng mô hình tĩnh của DeepFace không thể tự thích ứng với thuộc tính này.
\end{itemize}

\subsection{Mục tiêu nghiên cứu}
Đề tài hướng tới việc cải tiến thuật toán của thư viện DeepFace để tối ưu hóa độ chính xác so khớp khuôn mặt trước và sau phẫu thuật thẩm mỹ. Các mục tiêu cụ thể bao gồm:
\begin{enumerate}
    \item Triển khai giải thuật \textbf{Tích hợp vùng cục bộ thích ứng (Adaptive Local Patch Fusion)} can thiệp trực tiếp vào quy trình so khớp của DeepFace: chia khuôn mặt thành 3 vùng cục bộ độc lập (Mắt/Trán, Mũi, Miệng/Cằm).
    \item Phát triển thuật toán \textbf{Điều chỉnh trọng số động (Dynamic Weighting Adjustment)} để tự động nhận biết vùng bị phẫu thuật thông qua độ lệch khoảng cách Cosine, từ đó giảm trọng số của vùng phẫu thuật và tăng trọng số của vùng bất biến.
    \item \textbf{Xây dựng hệ thống phần mềm hoàn chỉnh bao gồm Go Backend và Next.js Frontend để demo trực quan thời gian thực cơ chế khoanh vùng, tính điểm tương đồng cục bộ và đưa ra quyết định so khớp thích ứng.}
    \item Đánh giá thực nghiệm hiệu năng của giải pháp đề xuất trên tập ảnh kiểm thử trước/sau phẫu thuật thẩm mỹ để chứng minh sự vượt trội về độ chính xác so với giải thuật DeepFace gốc.
\end{enumerate}

\subsection{Phạm vi nghiên cứu}
\begin{itemize}
    \item \textbf{Phạm vi công nghệ:} Tập trung cải tiến mã nguồn trích xuất đặc trưng của thư viện DeepFace chạy trên mô hình VGG-Face. Tích hợp giải thuật cắt phân vùng hình học cục bộ không chạy lại bộ dò để đảm bảo tốc độ xử lý nhanh.
    \item \textbf{Dữ liệu kiểm thử:} Sử dụng các tập ảnh khuôn mặt trước/sau các ca phẫu thuật thẩm mỹ phổ biến (nâng mũi, gọt cằm, cắt mí).
\end{itemize}

\subsection{Bố cục đồ án}
\begin{itemize}
    \item \textbf{Chương 2: Cơ sở lý thuyết và Tổng quan tài liệu} trình bày về công nghệ nhận diện khuôn mặt, các mô hình CNN trong DeepFace, và các nghiên cứu liên quan đến bài toán nhận diện sau phẫu thuật thẩm mỹ.
    \item \textbf{Chương 3: Phương pháp nghiên cứu và Thiết kế hệ thống} chi tiết hóa giải thuật đề xuất về cắt phân vùng cục bộ, tính khoảng cách Cosine cục bộ và giải thuật phân bổ trọng số thích ứng động.
    \item \textbf{Chương 4: Kết quả thực nghiệm và Thảo luận} trình bày về kịch bản kiểm thử, so sánh chỉ số khoảng cách và phân tích tính hiệu quả của mô hình đề xuất trên giao diện ứng dụng.
    \item \textbf{Chương 5: Kết luận và Hướng phát triển} tổng kết những đóng góp của đồ án và hướng cải tiến trong tương lai.
\end{itemize}
